\bookchapter{付録F：折りたみ画面の選び方と使い方特集}

\chapterintro{—— あなたのデジタルライフガイド}

\booksection{一、折りたたみ画面は今どこまで発展しているか？}

折りたたみ画面は「未来のテクノロジー」から「今買える」成熟カテゴリーになりました。しかし、ストレート型スマートフォンとはフォームファクター、使用習慣、メンテナンス方法に大きな違いがあり、買う価値があるか、どう選ぶか、どう使うか、この記事で整理しましょう。\textbf{一、折りたたみ画面は今どこまで発展しているか？}先に結論を言うと：

折りたたみ画面はもはや「お試し玩具」ではありませんが、「誰もが買うべき」主流のフォームファクターでもありません

\begin{itemize}
  \item \textbf{。ニッチなギーク向け製品と大衆消費財の間に位置し、特定のユーザー層に適しています。}現在主流は2つのカテゴリーに分かれます：
  \item \textbf{横折り（ブック型）}：閉じると普通のスマートフォンサイズ、開くとほぼ正方形の小型タブレット。代表：Huawei Mate Xシリーズ、Samsung Galaxy Z Foldシリーズ、Xiaomi MIX Foldシリーズ、vivo X Foldシリーズ。「1台のスマートフォン＝スマートフォン＋タブレット」が売り。
\end{itemize}

2024-2026年の主要な進歩：ヒンジ寿命が普遍的に50万回の折りたたみを保証（1日100回開閉で13年使用可）、UTG超薄ガラスが初期のプラスチックフィルムに取って代わり、折り目は目視で確認できるが手触りはすでに非常に軽微、価格が1万元帯から4000-5000元帯に下下探。

\begin{tcbinfo}{：閉じるとパウダーケースのよう、開くと通常サイズのスマートフォン。代表：Huawei Pocketシリーズ、Samsung Galaxy Z Flipシリーズ、OPPO Find N Flip、Xiaomi MIX Flip。「小型携帯＋外画面クイックチェック」が売り。}
2024-2025年の重要な進歩：ヒンジ寿命が一般的に50万回の折りたたみに対応（1日100回開閉で13年使用可能）、UTG超薄ガラスが初期のプラスチックフィルムに代わり、折り目は肉眼で見えるが触感は非常に軽微、価格は1万元帯から5000-6000元帯に下落。

\end{tcbinfo}

\booksection{二、どんな人が折りたたみスマートフォンに向いているか？}

現在ストレート型スマートフォンで完全に十分で、「画面が足りない」というペインポイントがないなら、折りたたみスマートフォンは買わないでください。それは「錦上花を添える」デバイスであり、「雪中に炭を送る」デバイスではありません。

\begin{enumerate}
  \item \textbf{二、どんな人が折りたたみスマートフォンに向いているか？}実際のユーザーフィードバックに基づくと、以下のタイプの人が折りたたみスマートフォンを買って「マジで良かった」と感じる確率が高いです：
  \item \textbf{ヘビーリーダー/ドキュメント処理者}：スマートフォンでよくPDFを見たり、ドキュメントに注釈を付けたり、長文を読んだりする人。展開後の約7.5インチの内画面なら、A4 PDFをズームなしで見られ、体験が劇的に変わります。
  \item \textbf{頻繁に出張/会議のあるビジネスユーザー}：タブレットを持ちたくないがスマートフォンの画面が小さいと感じる人。横折りスマートフォンを展開してPPTを見たり、契約書にサインしたり、2画面でメールを照合したりでき、ストレート型より一段と効率的。
  \item \textbf{「ハングモード」でのクリエイティブ作業が好きな人}：90度に折りたたんで机に置き、上半分でファインダー下半分でシャッター（自撮り/低速度撮影）、上半分で動画下半分でメモ、上半分でライブ配信下半分でコメント確認——これはストレート型にはできないフォームファクターです。
\end{enumerate}

\begin{tcbwarn}{携帯性と大画面を両立したい女性ユーザー}
\begin{itemize}
  \item ：縦折りスマートフォンは閉じると小さなバッグに入り、開くと通常のスマートフォン、外画面は化粧鏡や通知確認にも使えます。
  \item こういう人は買わないで
  \item スマートフォンをショート動画視聴、WeChat、電話にしか使わない人——折りたたみ画面は何のメリットもなく、むしろより重く厚くなります。
\end{itemize}

\end{tcbwarn}

\booksection{三、選び方のポイント：この6つの指標を見る}

予算が厳しく、1台のスマートフォンを4年以上使いたい人——折りたたみスマートフォンは中古の価値保持率がストレート型より低く、修理コストも高いです。

\begin{itemize}
  \item \textbf{三、選び方のポイント：この6つの指標を見る}「折りたたみ」という言葉に目を奪われないで、選ぶときはここを重点的に見て：
  \item \textbf{ヒンジの種類}：水滴ヒンジ（折り目が浅く、寿命が長い、現在主流）がU字ヒンジ（折り目が深い、初期の設計）より優れています。店員に聞くかレビューで確認。
  \item \textbf{内画面の素材}：実店舗で自分の手で展開して見てください。2026年の合格ラインは：画面点灯時は正面からほとんど見えず、斜めの光で浅い溝が見え、指で撫でると軽い凹みがあるが手を切らない。
  \item \textbf{折り目の状態}：実店舗で実際に展開して見て。2025年の合格ラインは：点灯時は正面からはほとんど見えず、斜めの光で浅い溝が見え、指でなぞるとわずかなへこみがあるが手を切りません。
  \item \textbf{外画面サイズ}：横折りスマートフォンは外画面が大きいほど良い（6.4インチ以上が理想）、閉じた状態で通常のスマートフォンとして使える。縦折りスマートフォンは外画面が大きいほど実用的（4インチ以上ならメッセージ返信、ナビ確認が可能）。
  \item \textbf{重量と厚さ}：横折り展開後230-260gが一般的、230g未満は優秀。閉じた厚さ12-14mm、12mm未満は薄い。縦折りは閉じた状態でシャツのポケットに楽に入れられるのが合格ライン。
\end{itemize}

\begin{tcbpractice}{アフターと保証}
\begin{enumerate}
  \item ：内画面割れの修理価格を重点的に確認。一部ブランドは初年度1回無料で内画面交換、自己負担の場合は2000-4000元かかることも。これは折りたたみスマートフォン最大の隠れコストです。
  \item 実店舗体験チェックリスト
  \item 実際に10回開閉し、ヒンジのダンピングが滑らかか、異音がないか確認。
  \item 展開後、ライトに斜めに向けて折り目を見て、それから正面から点灯画面を見る。
  \item 展開状態を片手で30秒持ち、重量が許容できるか確認。
\end{enumerate}

\end{tcbpractice}

\booksection{四、折りたたみスマートフォンならではの使い方のコツ}

内画面単独の修理価格と保証ポリシーを確認。

\booksubsection{四、折りたたみスマートフォンならではの使い方のコツ}

買った後、これらの使い方はストレート型では体験できません：

\booksubsection{分割画面とマルチタスク}

横折りスマートフォン展開後、ほとんどのシステムが上下/左右分割画面に対応。典型的シナリオ：上半分でオンライン講座の動画、下半分でノート、左にWeChat右にドキュメント、上にマップナビ下に音楽。ストレート型の分割画面よりはるかに実用的で、画面が十分大きく、どちらの半分も窮屈ではありません。

\begin{itemize}
  \item \textbf{2. ハングモード（90度折りたたみ）}これは折りたたみ画面の「切り札」的な使い方です：
  \item \textbf{自撮り/集合写真}：机に置いて90度に折り、上半分でファインダー、下半分でシャッター、人に撮ってもらう必要なし。
  \item \textbf{低速度撮影/夜景}：同じく机に置いて安定撮影、三脚が不要。
  \item \textbf{ビデオ通話}：両手がフリーになり、相手から見る角度もより自然。
  \item \textbf{動画視聴}：机に置いてミニテレビに、飛行機/新幹線のテーブルで特に便利。
\end{itemize}

\booksubsection{外画面プレビュー＋内画面撮影}

：背面カメラで自分を撮るとき、外画面にリアルタイムで映像が表示され、構図がより正確に。

\booksubsection{外画面クイックチェック}

縦折りスマートフォンを閉じた後、外画面で：時間/通知/天気の確認、WeChatへの返信（音声テキストまたはクイック返信）、交通系ICカードのタッチ、ナビの確認、化粧鏡（インカメラの起動）が可能。毎回展開する必要がなく、省電力で便利。

\booksection{五、メンテナンスのポイント：折りたたみ画面を長持ちさせる}

一部の横折りスマートフォン（Huawei、Samsung）は展開後「PCモード」に切り替え可能、外部モニターとキーボード·マウスを接続すれば簡易PCのように、一時的なドキュメント処理やメール返信には十分。メインの生産性ツールではありませんが、緊急時には非常に便利。

\begin{itemize}
  \item \textbf{五、メンテナンスのポイント：折りたたみ画面を長持ちさせる}折りたたみ画面で最もデリケートなのが内画面とヒンジ、これらの習慣で寿命を大幅に延ばせます：
  \item \textbf{低温では無理に開閉しない}：0℃以下ではUTGガラスが脆くなる。開閉する前に懐で数分温めて。これは公式マニュアルにも書いてあるが誰も読まない重要ポイント。
  \item \textbf{爪/硬いもので内画面に触れない}：内画面には保護層があるが外画面より柔らかい。指の腹で操作し、爪で突かない。
  \item \textbf{閉じるときに画面内に異物がないか確認}：一粒の砂が挟まったまま閉じると、内画面を直接圧壊する可能性。閉じる前に一目確認。
  \item \textbf{普通の強化ガラスフィルムは貼らない}：内画面には公式専用フィルムまたはUTG保護フィルムのみ使用。普通の強化ガラスフィルムは折りたたみとタッチ操作に影響。公式サービスセンターでのフィルム貼りをおすすめ。
  \item \textbf{ヒンジにほこりや液体を入れない}：IPX8防水を謳うものが多いが、防塵等級は一般的に低い。ビーチや工事現場などの環境では注意。
\end{itemize}

\textbf{開閉の力を均一に}\par

\begin{itemize}
  \item ：片手で勢いよく開かない、両手で中央から均等に展開するのが最も安定。
  \item 内画面破損の多発シーン
  \item 閉じるときにイヤホンケーブル/充電ケーブルを挟む——内画面が直接圧壊。
  \item バッグの中で鍵と一緒に入れる——外画面は無事でも、内画面が鍵でひび割れ。
\end{itemize}

\booksection{六、折りたたみスマートフォンの「落とし穴」と心理的期待値}

1年使って「折り目が深くなった」と感じて純正フィルムを剥がす——フィルムを剥がすと内画面も一緒に破損。

\begin{itemize}
  \item \textbf{六、折りたたみスマートフォンの「落とし穴」と心理的期待値}：2026年の技術では「完全に痕なしが」は実現できず、「画面点灯時正面からほとんど見えない」までしかできません。完璧主義者で、毎日折り目を見て不快に感じるなら、買わないでください。
  \item \textbf{折り目はずっとある}：2025年の技術では「完全に跡なし」は実現できず、「点灯時正面からはほとんど見えない」が限界。完璧主義者で、毎日折り目を見てつらいなら、買わないで。
  \item \textbf{より重くより厚い}：横折りスマートフォンは展開後一般的に230g以上で、片手持ちはストレート型より疲れる。閉じた厚さ12-14mmで、ズボンのポケットがより膨らむ。
  \item \textbf{バッテリー持ちは一般的に普通}：ヒンジ機構を収めるため、バッテリー空間が圧縮され、同価格帯のバッテリー持ちは通常ストレート型に劣る。
  \item \textbf{価値保持率が低い}：1年後の中古価格は通常新品の50-60\%で、ストレート型より下落が早い。
  \item \textbf{修理が高い}：内画面ユニットの修理価格2000-4000元、ヒンジ修理500-1500元。購入前に保証ポリシーを確認。
\end{itemize}

アプリの対応がまちまち

：主流アプリは大画面に対応済みだが、ニッチなアプリは単に「引き伸ばし拡大」の可能性があり、体験はいまひとつ。

\booksection{七、主流ブランドの選び方（2026年）}

—— 折りたたみスマートフォンの選び方の哲学について

\begin{itemize}
  \item \textbf{七、主要ブランドの選び方（2025年）}具体的な型番は書かず（古くなるため）、ブランドの特徴だけ書き、素早く絞り込めるようにします：
  \item \textbf{Huawei Mate Xシリーズ}：横折りのベンチマーク、HarmonyOSエコシステムとの連携が強力、外画面の比率が通常のスマートフォンに近く、閉じた状態でメイン機として使える。価格は高いが価値保持率は相対的に最も良い。
  \item \textbf{Samsung Galaxy Z Fold/Flip}：折りたたみ画面の元祖、ヒンジの調整が成熟、Z Flipは縦折りのルックス担当。中国国内のサービス拠点が少ないのが弱点。
  \item \textbf{Xiaomi MIX Fold/Flip}：コストパフォーマンス路線、スペックが厚く、価格は比較的お手頃。試してみたいがあまりお金をかけたくないユーザーに適している。
  \item \textbf{vivo X Foldシリーズ}：横折りの中で内外画面の比率が最も通常のスマートフォンに近く、閉じてストレート型として使う体験が良い。
  \item \textbf{OPPO Find Nシリーズ}：横折りの中でサイズがややコンパクト、片手で持てる、大型折りたたみが重いと感じるユーザーに適している。
\end{itemize}

Huawei Pocketシリーズ

\begin{tcbinfo}{最後の注意}
折りたたみ画面技術は毎年イテレーションしており、本付録は2026年に執筆。購入前に必ず最新のレビューとユーザーフィードバックを確認し、ヒンジ寿命の実測、折り目の表現、内画面の修理ポリシーの3項目に注目。実店舗で直接触れることがレビューを見るより常に重要です。

\end{tcbinfo}
